\documentclass[10pt,a4paper]{amsart}
\usepackage{amsmath,amssymb,amsthm,graphicx}
% Retained minimal preamble for the Stacks Project chapter more-morphisms.tex (More on Morphisms).
% Verbatim excerpt of the chapter preamble: Xy-pic setup (native preamble lines 14-18) and the native
% theorem declarations with the mathematical macros (native preamble lines 154-217). Copied mechanically;
% nothing added, removed or redefined. The wrapper class and the amsmath/amssymb/amsthm/graphicx base
% packages are supplied by the assembler, so the class/externaldocument part of the native preamble is omitted.
% For commutative diagrams we use Xy-pic
\usepackage[all]{xy}

% We use 2cell for 2-commutative diagrams.
\xyoption{2cell}
% Theorem environments.
%
\theoremstyle{plain}
\newtheorem{theorem}[subsection]{Theorem}
\newtheorem{proposition}[subsection]{Proposition}
\newtheorem{lemma}[subsection]{Lemma}

\theoremstyle{definition}
\newtheorem{definition}[subsection]{Definition}
\newtheorem{example}[subsection]{Example}
\newtheorem{exercise}[subsection]{Exercise}
\newtheorem{situation}[subsection]{Situation}

\theoremstyle{remark}
\newtheorem{remark}[subsection]{Remark}
\newtheorem{remarks}[subsection]{Remarks}

\numberwithin{equation}{subsection}

% Macros
%
\def\lim{\mathop{\mathrm{lim}}\nolimits}
\def\colim{\mathop{\mathrm{colim}}\nolimits}
\def\Spec{\mathop{\mathrm{Spec}}}
\def\Hom{\mathop{\mathrm{Hom}}\nolimits}
\def\Ext{\mathop{\mathrm{Ext}}\nolimits}
\def\SheafHom{\mathop{\mathcal{H}\!\mathit{om}}\nolimits}
\def\SheafExt{\mathop{\mathcal{E}\!\mathit{xt}}\nolimits}
\def\Sch{\mathit{Sch}}
\def\Mor{\mathop{\mathrm{Mor}}\nolimits}
\def\Ob{\mathop{\mathrm{Ob}}\nolimits}
\def\Sh{\mathop{\mathit{Sh}}\nolimits}
\def\NL{\mathop{N\!L}\nolimits}
\def\CH{\mathop{\mathrm{CH}}\nolimits}
\def\proetale{{pro\text{-}\acute{e}tale}}
\def\etale{{\acute{e}tale}}
\def\QCoh{\mathit{QCoh}}
\def\Ker{\mathop{\mathrm{Ker}}}
\def\Im{\mathop{\mathrm{Im}}}
\def\Coker{\mathop{\mathrm{Coker}}}
\def\Coim{\mathop{\mathrm{Coim}}}

% Boxtimes
%
\DeclareMathSymbol{\boxtimes}{\mathbin}{AMSa}{"02}

%
% Macros for moduli stacks/spaces
%
\def\QCohstack{\mathcal{QC}\!\mathit{oh}}
\def\Cohstack{\mathcal{C}\!\mathit{oh}}
\def\Spacesstack{\mathcal{S}\!\mathit{paces}}
\def\Quotfunctor{\mathrm{Quot}}
\def\Hilbfunctor{\mathrm{Hilb}}
\def\Curvesstack{\mathcal{C}\!\mathit{urves}}
\def\Polarizedstack{\mathcal{P}\!\mathit{olarized}}
\def\Complexesstack{\mathcal{C}\!\mathit{omplexes}}
% \Pic is the operator that assigns to X its picard group, usage \Pic(X)
% \Picardstack_{X/B} denotes the Picard stack of X over B
% \Picardfunctor_{X/B} denotes the Picard functor of X over B
\def\Pic{\mathop{\mathrm{Pic}}\nolimits}
\def\Picardstack{\mathcal{P}\!\mathit{ic}}
\def\Picardfunctor{\mathrm{Pic}}
\def\Deformationcategory{\mathcal{D}\!\mathit{ef}}
\usepackage{cleveref}
\crefname{originalreference}{original reference}{original references}
\Crefname{originalreference}{Original reference}{Original references}
\makeatletter
\newlabel{coherent-lemma-quasi-coherent-affine-cohomology-zero}{{01XB}{source}{}{sourcepointer-1}{}}
\newlabel{coherent-lemma-quasi-coherent-affine-cohomology-zero@cref}{{[originalreference][0][]01XB}{[1][0][]source}{}{}{}}
\makeatother
\crefname{originalreference}{original reference}{original references}
\Crefname{originalreference}{Original reference}{Original references}
\makeatletter
\newlabel{coherent-lemma-relative-affine-cohomology}{{089W}{source}{}{sourcepointer-2}{}}
\newlabel{coherent-lemma-relative-affine-cohomology@cref}{{[originalreference][0][]089W}{[1][0][]source}{}{}{}}
\makeatother
\crefname{originalreference}{original reference}{original references}
\Crefname{originalreference}{Original reference}{Original references}
\makeatletter
\newlabel{coherent-lemma-relative-affine-vanishing}{{01XC}{source}{}{sourcepointer-3}{}}
\newlabel{coherent-lemma-relative-affine-vanishing@cref}{{[originalreference][0][]01XC}{[1][0][]source}{}{}{}}
\makeatother
\crefname{originalreference}{original reference}{original references}
\Crefname{originalreference}{Original reference}{Original references}
\makeatletter
\newlabel{cohomology-lemma-quasi-separated-cohomology-colimit}{{01FF}{source}{}{sourcepointer-4}{}}
\newlabel{cohomology-lemma-quasi-separated-cohomology-colimit@cref}{{[originalreference][0][]01FF}{[1][0][]source}{}{}{}}
\makeatother
\crefname{originaldefinition}{original definition}{original definitions}
\Crefname{originaldefinition}{Original definition}{Original definitions}
\makeatletter
\newlabel{definition-affine-stratification-number}{{0F2V}{source}{}{sourcepointer-5}{}}
\newlabel{definition-affine-stratification-number@cref}{{[originaldefinition][0][]0F2V}{[1][0][]source}{}{}{}}
\makeatother
\crefname{originalreference}{original reference}{original references}
\Crefname{originalreference}{Original reference}{Original references}
\makeatletter
\newlabel{morphisms-section-support}{{056H}{source}{}{sourcepointer-6}{}}
\newlabel{morphisms-section-support@cref}{{[originalreference][0][]056H}{[1][0][]source}{}{}{}}
\makeatother
\crefname{originalreference}{original reference}{original references}
\Crefname{originalreference}{Original reference}{Original references}
\makeatletter
\newlabel{properties-lemma-quasi-coherent-colimit-finite-type}{{01PG}{source}{}{sourcepointer-7}{}}
\newlabel{properties-lemma-quasi-coherent-colimit-finite-type@cref}{{[originalreference][0][]01PG}{[1][0][]source}{}{}{}}
\makeatother
\crefname{originalreference}{original reference}{original references}
\Crefname{originalreference}{Original reference}{Original references}
\makeatletter
\newlabel{schemes-section-quasi-coherent}{{01LA}{source}{}{sourcepointer-8}{}}
\newlabel{schemes-section-quasi-coherent@cref}{{[originalreference][0][]01LA}{[1][0][]source}{}{}{}}
\makeatother
\title{Cohomology of quasi-coherent modules vanishes above the affine stratification number}
\author{The Stacks Project Authors}
\date{stacks-project:more-morphisms.tex:local-snapshot-20261011; GNU FDL 1.2 or later}
\begin{document}
\maketitle
\noindent\textit{Editorial note.} Counted span is the author's own proof environment, source lines 21645-21693, immediately following the statement at 21638-21643.\par\medskip
\noindent\textit{Editorial note.} Seven cross-chapter pointers (Cohomology of Schemes, Cohomology, Properties, Morphisms) plus the chapter's own definition are recorded in assembly.external\_references.\par\medskip
\noindent\textit{Editorial note.} The reduction to finite type modules supported on Z is displayed in the span before the support trick is applied.\par\medskip
\noindent\textit{Allowed standard prerequisites (editorial).}\begin{itemize}
\item Definition of the affine stratification number. chapter definition definition-affine-stratification-number, cited at proof line 21650. Supplies the inductive structure: an affine open U whose complement carries a stratification number one smaller.
\item Vanishing of quasi-coherent cohomology on an affine scheme. Cohomology of Schemes, lemma coherent-lemma-quasi-coherent-affine-cohomology-zero, cited at proof lines 21648-21649. Provides the base case n = 0.
\item Relative affine vanishing and affine cohomology comparison. Cohomology of Schemes, lemmas coherent-lemma-relative-affine-vanishing and coherent-lemma-relative-affine-cohomology, cited in the proof. Transfers the computation to the complement Z and identifies H\textasciicircum{}p(X, G) with H\textasciicircum{}p(Z', G).
\item Colimits of quasi-coherent modules of finite type. Cohomology, lemma cohomology-lemma-quasi-separated-cohomology-colimit and Properties, lemma properties-lemma-quasi-coherent-colimit-finite-type, cited in the proof. Reduces the general quasi-coherent F to finite type modules supported on Z.
\end{itemize}\medskip
\ifdefined\hypertarget\hypertarget{sourcepointer-1}{}\fi
\noindent\textit{Source reference.} Original reference 01XB: focused-seeds/geometry/S9-morph-vanishing-affine-strat-number/ref-tags.txt, source lines 4--4. Navigation pointer only: the counted proof cites this statement by \textbackslash{}ref as a standing input; the deduction the counted proof has to make is carried out inside the counted span, and no part of that deduction is delegated to this pointer.\par\medskip
\ifdefined\hypertarget\hypertarget{sourcepointer-2}{}\fi
\noindent\textit{Source reference.} Original reference 089W: focused-seeds/geometry/S9-morph-vanishing-affine-strat-number/ref-tags.txt, source lines 5--5. Navigation pointer only: the counted proof cites this statement by \textbackslash{}ref as a standing input; the deduction the counted proof has to make is carried out inside the counted span, and no part of that deduction is delegated to this pointer.\par\medskip
\ifdefined\hypertarget\hypertarget{sourcepointer-3}{}\fi
\noindent\textit{Source reference.} Original reference 01XC: focused-seeds/geometry/S9-morph-vanishing-affine-strat-number/ref-tags.txt, source lines 6--6. Navigation pointer only: the counted proof cites this statement by \textbackslash{}ref as a standing input; the deduction the counted proof has to make is carried out inside the counted span, and no part of that deduction is delegated to this pointer.\par\medskip
\ifdefined\hypertarget\hypertarget{sourcepointer-4}{}\fi
\noindent\textit{Source reference.} Original reference 01FF: focused-seeds/geometry/S9-morph-vanishing-affine-strat-number/ref-tags.txt, source lines 7--7. Navigation pointer only: the counted proof cites this statement by \textbackslash{}ref as a standing input; the deduction the counted proof has to make is carried out inside the counted span, and no part of that deduction is delegated to this pointer.\par\medskip
\ifdefined\hypertarget\hypertarget{sourcepointer-5}{}\fi
\noindent\textit{Source reference.} Original definition 0F2V: sources/stacks-more-morphisms/pinned-original-source.tex, source lines 21569--21569. Navigation pointer only: the counted proof cites this statement by \textbackslash{}ref as a standing input; the deduction the counted proof has to make is carried out inside the counted span, and no part of that deduction is delegated to this pointer.\par\medskip
\ifdefined\hypertarget\hypertarget{sourcepointer-6}{}\fi
\noindent\textit{Source reference.} Original reference 056H: focused-seeds/geometry/S9-morph-vanishing-affine-strat-number/ref-tags.txt, source lines 8--8. Navigation pointer only: the counted proof cites this statement by \textbackslash{}ref as a standing input; the deduction the counted proof has to make is carried out inside the counted span, and no part of that deduction is delegated to this pointer.\par\medskip
\ifdefined\hypertarget\hypertarget{sourcepointer-7}{}\fi
\noindent\textit{Source reference.} Original reference 01PG: focused-seeds/geometry/S9-morph-vanishing-affine-strat-number/ref-tags.txt, source lines 9--9. Navigation pointer only: the counted proof cites this statement by \textbackslash{}ref as a standing input; the deduction the counted proof has to make is carried out inside the counted span, and no part of that deduction is delegated to this pointer.\par\medskip
\ifdefined\hypertarget\hypertarget{sourcepointer-8}{}\fi
\noindent\textit{Source reference.} Original reference 01LA: focused-seeds/geometry/S9-morph-vanishing-affine-strat-number/ref-tags.txt, source lines 10--10. Navigation pointer only: the counted proof cites this statement by \textbackslash{}ref as a standing input; the deduction the counted proof has to make is carried out inside the counted span, and no part of that deduction is delegated to this pointer.\par\medskip
\begin{proposition}
\label{proposition-vanishing-affine-stratification-number}
Let $X$ be a nonempty quasi-compact and quasi-separated scheme with
affine stratification number $n$. Then $H^p(X, \mathcal{F}) = 0$, $p > n$
for every quasi-coherent $\mathcal{O}_X$-module $\mathcal{F}$.
\end{proposition}

\begin{proof}
We will prove this by induction on the affine stratification number $n$.
If $n = 0$, then $X$ is affine and the result is
Cohomology of Schemes, Lemma
\ref{coherent-lemma-quasi-coherent-affine-cohomology-zero}.
Assume $n > 0$. By Definition \ref{definition-affine-stratification-number}
there is an affine scheme $U$ and an affine open immersion $j : U \to X$
such that the complement $Z$ has affine stratification number $n - 1$.
As $U$ and $j$ are affine we have $H^p(X, j_*(\mathcal{F}|_U)) = 0$
for $p > 0$, see Cohomology of Schemes, Lemmas
\ref{coherent-lemma-relative-affine-cohomology} and
\ref{coherent-lemma-relative-affine-vanishing}.
Denote $\mathcal{K}$ and $\mathcal{Q}$ the kernel and cokernel
of the map $\mathcal{F} \to j_*(\mathcal{F}|_U)$. Thus we obtain
an exact sequence
$$
0 \to \mathcal{K} \to \mathcal{F} \to j_*(\mathcal{F}|_U)
\to \mathcal{Q} \to 0
$$
of quasi-coherent $\mathcal{O}_X$-modules (see
Schemes, Section \ref{schemes-section-quasi-coherent}).
A standard argument, breaking our exact sequence into short
exact sequences and using the long exact cohomology sequence,
shows it suffices to prove $H^p(X, \mathcal{K}) = 0$
and $H^p(X, \mathcal{Q}) = 0$ for $p \geq n$.
Since $\mathcal{F} \to j_*(\mathcal{F}|_U)$ restricts
to an isomorphism over $U$, we see that $\mathcal{K}$
and $\mathcal{Q}$ are supported on $Z$.
By Properties, Lemma
\ref{properties-lemma-quasi-coherent-colimit-finite-type}
we can write these modules as the filtered colimits of
their finite type quasi-coherent submodules.
Using the fact that cohomology of sheaves on $X$ commutes
with filtered colimits, see
Cohomology, Lemma \ref{cohomology-lemma-quasi-separated-cohomology-colimit},
we conclude it suffices to show that if $\mathcal{G}$
is a finite type quasi-coherent module whose support
is contained in $Z$, then $H^p(X, \mathcal{G}) = 0$ for $p \geq n$.
Let $Z' \subset X$ be the scheme theoretic support of
$\mathcal{G} \oplus \mathcal{O}_Z$; we may and do think
of $\mathcal{G}$ as a quasi-coherent module on $Z'$, see
Morphisms, Section \ref{morphisms-section-support}.
Then $Z'$ and $Z$ have the same underlying topological space
and hence the same affine stratification number, namely $n - 1$.
Hence $H^p(X, \mathcal{G}) = H^p(Z', \mathcal{G})$
(equality by Cohomology of Schemes, Lemma
\ref{coherent-lemma-relative-affine-cohomology})
vanishes for $p \geq n$ by induction hypothesis.
\end{proof}

\end{document}
